% (c) 2026 Stephanie Märklin, OST Ostschweizer Fachhochschule
%
% !TEX root = ../../buch.tex
% !TEX encoding = UTF-8
%
% Rolle im Paper: Die Frage in der Sprache des Modells exakt formulieren —
% Übergabe an die Mathematik
%
% Stichworte:
% 17.3 Stabilität als mathematische Eigenschaft
% - Andocken: "Es bleibt, den Stabilitätsbegriff zu übersetzen."
% - Übersetzung als Kette: Ruhelage -> x = 0 (Begründung mit denn),
%   Störung -> x(0) != 0, Rückkehr -> x(t) gegen 0
% - Definition (heisst-Muster): stabil, wenn jede Lösung für jeden
%   Anfangszustand gegen 0 geht
% - Schluss-Feststellung: Flugzeugfrage -> DGL-Frage; Suchauftrag
%   "Wir suchen x(t), das die Gleichung erfüllt" (Whiteboard/Vortrag);
%   17.4 dockt mit "Ableitung proportional zur Funktion" an
% - Stil: genau EIN Doppelpunkt (vor dem Suchauftrag); keine Fragen

\section{Stabilität als mathematische Eigenschaft
\label{aircraft:section:3-mathematische-stabilitaet}}
\kopfrechts{Mathematische Stabilität}
Wir übertragen nun auch den Stabilitätsbegriff aus
Abschnitt~\ref{aircraft:section:1-grundbegriffe} auf das dynamische System.
In der Ruhelage liegt die Querachse horizontal, die Längsachse zeigt in die
Bewegungsrichtung, und das Flugzeug dreht sich weder um die Längsachse noch
um die Hochachse.
Alle vier Zustandsgrössen der Seitenbewegung sind damit null.
In der Ruhelage ist also $\vec{x} = \vec{0}$.
Eine Störung ist ein Anfangszustand $\vec{x}(0) \neq \vec{0}$.
Die Rückkehr in die Ruhelage bedeutet, dass $\vec{x}(t)$ für
wachsende Zeit gegen $\vec{0}$ strebt.
Das System heisst also stabil, wenn jede Lösung der
Gleichung~\eqref{aircraft:eqn:bewegungsgesetz} für jeden
Anfangszustand gegen $\vec{0}$ geht.

Aus der Frage nach der Stabilität eines Flugzeugs ist damit eine
Frage über die Lösungen einer Differentialgleichung geworden.
Wir suchen die Funktionen $\vec{x}(t)$, welche die
Gleichung~\eqref{aircraft:eqn:bewegungsgesetz} erfüllen.

